\subsubsection{数据预处理}

问题四在问题二数据基础上加入附件4的实时电价序列。附件4按日、按10分钟记录价格，预处理先刻画价格特征，再构造可比的日前预测。

第一步是价格序列的对齐与统计。附件4的记录间隔同样为10分钟，与附件2的负载、光伏序列在时间上完全对齐，可直接按日、按时段组织为$p_t^{(d)}$。我们对其作统计刻画，得到全年实时电价的均值、标准差与变异系数：
\begin{equation}\label{eq:src47}
    \left\{\begin{aligned}
    \bar{p}&=\frac{1}{N}\sum_{d,t}p_{t}^{(d)}=0.7662,\\[1mm]
    \sigma_{p}&=\sqrt{\frac{1}{N}\sum_{d,t}\left(p_{t}^{(d)}-\bar{p}\right)^{2}}=0.3407,\\[1mm]
    \mathrm{CV}_{p}&=\frac{\sigma_{p}}{\bar{p}}=44.46\%
    \end{aligned}\right.
\end{equation}
其中，$N=365\times144$为全年样本总数。实时电价的最小值为0.0076元/kWh、最大值为1.7936元/kWh，峰谷比达236倍，而附件1的固定电价峰谷比仅3.76倍，说明实时电价的波动强度远高于固定电价。

第二步是价格不确定性度量。除全年的整体波动外，我们还需刻画电价在日内与逐日两个尺度上的波动强度：
\begin{equation}\label{eq:src48}
    \left\{\begin{aligned}
    \sigma_{\text{日内}}&=\frac{1}{365}\sum_{d}\mathrm{std}_{t}\left(p_{t}^{(d)}\right)=0.3229,\\[1mm]
    \sigma_{\text{逐日}}&=\mathrm{std}_{d}\left(\frac{1}{144}\sum_{t}p_{t}^{(d)}\right)=0.1036
    \end{aligned}\right.
\end{equation}
其中，$\sigma_{\text{日内}}=0.3229$元/kWh为365天日内电价标准差的均值，约为全年平均电价的42\%，说明日内价格差明显，储能可利用的价格差需要在效率与功率约束下判断；日均电价的取值范围为0.5378至0.9290元/kWh。这些指标描述实现价格的波动，不等同于预测误差或不可预测性。

第三步是价格预测的构造。与负载、光伏的预测同构，电价预测也采用“同时段基准＋偏差校正”的形式，且同样遵守因果性约束，即日前不使用当天未来实际价格；日内更新仅使用已完成时段的价格：
\begin{equation}\label{eq:src49}
    \hat{p}_{t}^{(d)}=\max\left\{p_{t}^{(d-7)}+\mathrm{median}_{k\in\mathcal K_d}\!\left(p_{t}^{(k)}-p_{t}^{(k-7)}\right),\ 10^{-4}\right\}
\end{equation}
其中，$p_t^{(d-7)}$为上周同日同时段的价格基准，$\mathcal K_d$为第$d$天之前、周同期数据已存在的最近最多$W=28$个已完成日。第二项为这些日期同一时段“周同比变化量”的中位数，起偏差校正作用；样本不足28日时使用已有样本，空集校正量取0。年初前7天以附件1电价作基准。下界$10^{-4}$用于避免价格为零时目标函数退化。式\eqref{eq:src49}的结构与式\eqref{eq:src18}一致，使得问题四可以复用问题二与问题三的求解流程。

经上述处理，我们得到实时价格的统计特征和日前预测序列，购电决策沿用问题二、问题三的储能与滚动调整框架。



\paragraph{季节性价格预测的统计解释}
式\eqref{eq:src49}以周周期季节性朴素预测为基准\cite{hyndman2021fpp}，并对最近最多28日同槽周同比变化作中位数校正。周周期项刻画同一星期位置的价格形态，中位数校正项刻画近期价格水平偏移，估计准则与式\eqref{eq:median_fit}一致。我们每天0时生成规划窗口的价格曲线，冻结当天24小时预测用于回测；日内仅在6、12、18时，以最近已完成时段的预测残差对当前价格曲线作衰减修正：
\begin{equation}\label{eq:price_update}
 \hat p_{t+j\mid t}=\max\{10^{-4},\hat p_{t+j}+\beta^j(p_t-\hat p_t)\},
 \qquad j\ge1,\quad\beta=0.90.
\end{equation}
其中$t$为最近已完成时段，$\beta$控制局部偏差的持续时间。日内校正与日前校正共享历史可用性约束，日前冻结值用于统一比较不同日期的预测误差。
